\documentclass[11pt]{article}
\usepackage[a4paper,margin=22mm]{geometry}
\usepackage{fontspec}
\setmainfont{DejaVu Serif}
\setsansfont{DejaVu Sans}
\usepackage{microtype}
\usepackage{xcolor}
\usepackage{hyperref}
\usepackage{enumitem}
\usepackage{parskip}
\definecolor{Accent}{HTML}{971D32}
\definecolor{Muted}{HTML}{666666}
\hypersetup{colorlinks=true,urlcolor=Accent,linkcolor=Accent}
\setlist{leftmargin=1.6em,itemsep=0.3em,topsep=0.35em}
\setcounter{tocdepth}{2}
\renewcommand{\contentsname}{Contents}
\newcommand{\sectionrule}{\par\vspace{-0.5em}\noindent\textcolor{Accent}{\rule{\linewidth}{0.6pt}}\par}
\begin{document}

\begin{center}
{\sffamily\bfseries\color{Accent} TOEFL WORD OF THE DAY\par}
\vspace{8mm}
{\fontsize{42}{48}\selectfont\bfseries globalization\par}
\vspace{2mm}
{\Large\sffamily\color{Accent} /ˌɡloʊbələˈzeɪʃən/\par}
\vspace{2mm}
{\small\sffamily Local audio phrase: globalization\par}
{\small\sffamily Audio file: audios/globalization.mp3\par}
\vspace{2mm}
{\large\sffamily\color{Muted} uncountable noun\par}
\vspace{5mm}
{\sffamily growing worldwide interconnectedness\par}
\vspace{6mm}
{\large Globalization is the process by which countries and societies become more closely connected through trade, technology, communication, culture, and migration.\par}
\vspace{6mm}
\begin{minipage}{0.84\textwidth}
\itshape “Advances in transportation and communication have accelerated globalization.”
\end{minipage}\par
\vspace{10mm}
\textbf{Date:} August 1st 2026\quad\textbullet\quad \textbf{Level:} B1--mid-B2 TOEFL
\vfill
Created and compiled by \textbf{Amir Kooshky}\par
\vspace{2mm}
\href{https://t.me/KooshkyTOEFL}{Telegram Channel}\quad\textbullet\quad
\href{https://www.instagram.com/kooshkytoefl}{Instagram}\quad\textbullet\quad
\href{https://t.me/KooshkyTOEFL\_pv}{Personal Telegram}
\end{center}
\newpage
\tableofcontents
\newpage

\section{Meanings and usage}
\sectionrule
\subsection{Meaning 1: Growing worldwide interconnectedness}
\textbf{Part of speech:} uncountable noun

\textbf{IPA:} /ˌɡloʊbələˈzeɪʃən/

\textbf{Pronunciation phrase:} globalization

\textbf{Local audio file:} audios/globalization.mp3

\textbf{Register:} neutral to formal; very common in economics, politics, sociology, business, history, and environmental studies

\textbf{Definition:} the process through which countries, economies, cultures, and people become increasingly connected and dependent on one another around the world

\textbf{In simpler words:} the world becoming more connected

\textbf{Typical pattern:} globalization of + industry, trade, culture, or economy

\subsubsection{Natural examples}
\begin{enumerate}
  \item Globalization has allowed companies to sell their products in markets around the world.
  \item Advances in transportation and communication have accelerated globalization.
  \item Critics argue that globalization can weaken local industries and traditions.
  \item The course examines how globalization has influenced employment and migration.
  \item Globalization has made it easier for ideas and cultural products to cross national borders.
  \item Some developing countries have benefited from globalization through increased trade and investment.
  \item The environmental effects of globalization extend far beyond individual national borders.
\end{enumerate}

\subsubsection{Nuance and implication}
\textbf{Central nuance:} The word describes a broad process involving the movement of goods, money, information, technology, people, and cultural influences across national borders.

\textbf{Usual implication:} Globalization may create opportunities and greater cooperation, but it may also increase inequality, economic dependence, cultural pressure, or environmental harm.

\textbf{Compared with internationalization:} Internationalization often describes a deliberate effort by a company, university, or organization to operate in multiple countries. Globalization is the broader process through which the world as a whole becomes increasingly interconnected.

\subsubsection{Grammar notes}
\textbf{How to use it:} Use globalization as the name of a general process.

\textbf{What can follow it:} It is often followed by of and a field or activity, as in the globalization of trade. It can also be followed by verbs such as increase, accelerate, transform, affect, or contribute.

\textbf{Common preposition:} Use the globalization of something when identifying the area becoming more global, as in the globalization of higher education.

\textbf{Noun use:} Globalization is normally uncountable, so do not use a or an before it. The form globalizations is extremely rare and is not normally needed by learners.

\textbf{Position in a sentence:} It can appear without the when discussing the process generally. Use the when referring to a specific form or aspect, as in the globalization of financial markets.

\subsubsection{Useful patterns}
\textbf{Pattern:} the globalization of + noun

\textbf{Use:} Use this to identify the industry, activity, institution, or cultural area becoming more internationally connected.

\textbf{Example:} The globalization of food production has created complex international supply chains.

\textbf{Pattern:} globalization has led to + noun

\textbf{Use:} Use this to describe a result of growing international interconnectedness.

\textbf{Example:} Globalization has led to greater competition between manufacturers.

\textbf{Pattern:} globalization has made it easier to + verb

\textbf{Use:} Use this to explain an activity that has become easier because national borders matter less.

\textbf{Example:} Globalization has made it easier to exchange scientific knowledge.

\textbf{Pattern:} the effects or impact of globalization

\textbf{Use:} Use this when discussing the consequences of globalization.

\textbf{Example:} The study investigates the impact of globalization on small farming communities.

\textbf{Pattern:} accelerate or drive globalization

\textbf{Use:} Use this for a force that causes countries and societies to become connected more rapidly.

\textbf{Example:} Digital technology has helped drive globalization.

\subsubsection{Common wrong usage}
\paragraph{Correction 1}
\textbf{Wrong:} The globalization has changed modern business.

\textbf{Correct:} Globalization has changed modern business.

\textbf{Why:} Do not normally use the when discussing globalization as a general process.

\paragraph{Correction 2}
\textbf{Wrong:} A globalization has created new economic opportunities.

\textbf{Correct:} Globalization has created new economic opportunities.

\textbf{Why:} Globalization is normally uncountable, so it is not used with a.

\paragraph{Correction 3}
\textbf{Wrong:} Globalization between trade has increased rapidly.

\textbf{Correct:} The globalization of trade has increased rapidly.

\textbf{Why:} Use the globalization of something, not globalization between something.

\paragraph{Correction 4}
\textbf{Wrong:} Globalization only means that companies sell products abroad.

\textbf{Correct:} Globalization includes international trade as well as growing cultural, technological, political, and social connections.

\textbf{Why:} The meaning is broader than business or trade alone.

\section{Word family}
\sectionrule
\subsection{global}
\textbf{Part of speech:} adjective

\textbf{IPA:} /ˈɡloʊbəl/

\textbf{Definition:} relating to the whole world or involving many parts of the world

\textbf{Relationship to the headword:} This adjective describes something involving or affecting the world as a whole.

\textbf{Grammar pattern:} global + economy, market, problem, network, trade, or population

\subsubsection{Examples}
\begin{enumerate}
  \item Climate change is a global problem that requires international cooperation.
  \item The company has developed a global network of suppliers.
\end{enumerate}

\subsection{globally}
\textbf{Part of speech:} adverb

\textbf{IPA:} /ˈɡloʊbəli/

\textbf{Definition:} throughout the world or in a way that involves the whole world

\textbf{Relationship to the headword:} This adverb describes something happening or existing worldwide.

\textbf{Grammar pattern:} operate, compete, spread, or be recognized globally

\subsubsection{Examples}
\begin{enumerate}
  \item The disease affects millions of people globally.
  \item The organization operates both locally and globally.
\end{enumerate}

\subsection{globalize}
\textbf{Part of speech:} transitive verb

\textbf{IPA:} /ˈɡloʊbəˌlaɪz/

\textbf{Definition:} to make a business, industry, activity, or system operate or develop throughout the world

\textbf{Relationship to the headword:} This verb describes causing something to become more global.

\textbf{Grammar pattern:} globalize + industry, market, business, production, or culture

\textbf{Usage note:} The verb is less common than globalization and is mainly used in formal economic and business contexts.

\subsubsection{Examples}
\begin{enumerate}
  \item The company plans to globalize its production network.
  \item Digital platforms have helped globalize the entertainment industry.
\end{enumerate}

\section{Common collocations}
\sectionrule
\subsection{economic globalization}
\textbf{Applies to:} Growing worldwide interconnectedness

\textbf{Meaning:} the increasing international connection of economies, markets, and businesses

\textbf{Example:} Economic globalization has increased trade and investment between countries.

\subsection{cultural globalization}
\textbf{Applies to:} Growing worldwide interconnectedness

\textbf{Meaning:} the spread and exchange of cultural products, practices, and ideas across countries

\textbf{Example:} Cultural globalization has helped music and films reach international audiences.

\subsection{the globalization of trade}
\textbf{Applies to:} Growing worldwide interconnectedness

\textbf{Meaning:} the process through which trade becomes increasingly international

\textbf{Pattern:} the globalization of + trade, industry, or market

\textbf{Example:} The globalization of trade has created highly complex supply chains.

\subsection{the process of globalization}
\textbf{Applies to:} Growing worldwide interconnectedness

\textbf{Meaning:} the gradual development of stronger connections between countries

\textbf{Example:} The process of globalization began long before the development of the internet.

\subsection{the forces of globalization}
\textbf{Applies to:} Growing worldwide interconnectedness

\textbf{Meaning:} the economic, technological, political, and cultural pressures that increase global connections

\textbf{Example:} Local businesses have struggled to adapt to the forces of globalization.

\subsection{the impact of globalization}
\textbf{Applies to:} Growing worldwide interconnectedness

\textbf{Meaning:} the effects that globalization has on people, institutions, or societies

\textbf{Pattern:} the impact of globalization on + group or area

\textbf{Example:} The article explores the impact of globalization on traditional communities.

\subsection{the effects of globalization}
\textbf{Applies to:} Growing worldwide interconnectedness

\textbf{Meaning:} the positive or negative consequences of growing global connections

\textbf{Pattern:} the effects of globalization on + noun

\textbf{Example:} The effects of globalization are not distributed equally among countries.

\subsection{the benefits of globalization}
\textbf{Applies to:} Growing worldwide interconnectedness

\textbf{Meaning:} the advantages created by stronger international connections

\textbf{Example:} Supporters emphasize the benefits of globalization for trade and innovation.

\subsection{the costs of globalization}
\textbf{Applies to:} Growing worldwide interconnectedness

\textbf{Meaning:} the negative consequences or disadvantages of globalization

\textbf{Example:} The documentary examines the social and environmental costs of globalization.

\subsection{accelerate globalization}
\textbf{Applies to:} Growing worldwide interconnectedness

\textbf{Meaning:} cause globalization to develop more rapidly

\textbf{Example:} Faster communication technologies have accelerated globalization.

\subsection{drive globalization}
\textbf{Applies to:} Growing worldwide interconnectedness

\textbf{Meaning:} act as an important cause of globalization

\textbf{Example:} International investment and digital communication continue to drive globalization.

\subsection{globalization has led to}
\textbf{Applies to:} Growing worldwide interconnectedness

\textbf{Meaning:} globalization has produced a particular result

\textbf{Pattern:} globalization has led to + noun or -ing form

\textbf{Example:} Globalization has led to greater cultural exchange between many societies.

\section{Synonyms and antonyms}
\sectionrule
\subsection{Close synonyms}
\subsubsection{global integration}
\textbf{Part of speech:} uncountable noun phrase

\textbf{Applies to:} Growing worldwide interconnectedness

\textbf{Shared meaning:} Both expressions describe increasing connections between countries and societies.

\textbf{Difference:} Global integration emphasizes different economies, systems, or societies becoming connected parts of a larger whole. Globalization is broader and can also include the spread of culture, information, technology, and social practices.

\textbf{Example:} Regional trade agreements have encouraged global integration.

\subsubsection{internationalization}
\textbf{Part of speech:} uncountable noun

\textbf{Applies to:} Growing worldwide interconnectedness

\textbf{Shared meaning:} Both involve increasing activity and connections across national borders.

\textbf{Difference:} Internationalization often describes a deliberate effort by a company, university, or institution to expand its international activities. Globalization usually refers to the wider historical process affecting the world as a whole.

\textbf{Example:} The university's internationalization strategy includes recruiting more overseas students.

\subsection{Direct or useful antonyms}
\subsubsection{deglobalization}
\textbf{Part of speech:} uncountable noun

\textbf{Applies to:} Growing worldwide interconnectedness

\textbf{Opposition:} Deglobalization is the process through which international economic, political, or social connections become weaker or less extensive.

\textbf{Usage note:} The word is especially common in economics and political analysis.

\textbf{Example:} Trade restrictions may contribute to deglobalization.

\section{Words learners may confuse}
\sectionrule
\subsection{globalism}
\textbf{Part of speech:} uncountable noun

\textbf{Why learners confuse them:} Both words relate to global connections, but one names a process while the other often names an idea or policy.

\textbf{Key difference:} Globalization is the process through which the world becomes more interconnected. Globalism is a belief, policy, or way of thinking that supports organizing activities and cooperation at a global level.

\textbf{globalization:} Globalization has increased the international movement of goods and information.

\textbf{globalism:} The politician criticized globalism and called for greater national independence.

\subsection{internationalization}
\textbf{Part of speech:} uncountable noun

\textbf{Why learners confuse them:} The two processes both involve activity across national borders and may occur together.

\textbf{Key difference:} Globalization is a broad worldwide process, while internationalization often refers to a specific organization deliberately expanding its international operations.

\textbf{globalization:} Globalization has transformed patterns of international trade.

\textbf{internationalization:} The internationalization of the university included creating exchange programs.

\section{Etymology}
\sectionrule
\textbf{Origin:} The word is formed from global and the noun-forming ending -ization, which expresses the process of making or becoming something.

\section{Practice exercises}
\sectionrule
\subsection{Word family choice}
Choose the best member of the word family for each blank.

\begin{enumerate}
  \item Climate change requires a coordinated \_\_\_\_\_ response.
  \item The company sells its products \_\_\_\_\_ through online platforms.
  \item Digital communication has contributed greatly to \_\_\_\_\_.
  \item The business hopes to \_\_\_\_\_ its services over the next decade.
\end{enumerate}

\subsection{Correct or incorrect?}
Decide whether each sentence is correct. Correct the inaccurate ones.

\begin{enumerate}
  \item Globalization has increased communication between distant regions.
  \item The globalization has created both opportunities and risks.
  \item The article discusses the globalization of higher education.
  \item A globalization has affected every national economy.
  \item Globalization and internationalization always mean exactly the same thing.
\end{enumerate}

\subsection{Collocation practice}
Complete each sentence with a natural collocation.

\begin{enumerate}
  \item The report examines the impact of globalization \_\_\_\_\_ rural communities.
  \item New transportation systems helped \_\_\_\_\_ globalization during the twentieth century.
  \item The globalization \_\_\_\_\_ trade has connected producers with distant markets.
  \item The course considers both the benefits and the \_\_\_\_\_ of globalization.
\end{enumerate}

\subsection{Complete answer key}
\subsubsection{Word family choice}
\begin{enumerate}
  \item \textbf{global.} Climate change requires a coordinated global response. \emph{Use the adjective global before the noun response.}
  \item \textbf{globally.} The company sells its products globally through online platforms. \emph{Use the adverb globally to describe where the company sells its products.}
  \item \textbf{globalization.} Digital communication has contributed greatly to globalization. \emph{Use the noun globalization after contributed to.}
  \item \textbf{globalize.} The business hopes to globalize its services over the next decade. \emph{Use the verb globalize after hopes to.}
\end{enumerate}

\subsubsection{Correct or incorrect?}
\begin{enumerate}
  \item \textbf{Correct} \emph{Globalization is correctly used without an article to refer to the general process.}
  \item \textbf{Incorrect}. Globalization has created both opportunities and risks. \emph{Do not normally use the when referring to globalization in general.}
  \item \textbf{Correct} \emph{The globalization of something is a standard pattern.}
  \item \textbf{Incorrect}. Globalization has affected every national economy. \emph{Globalization is normally uncountable and is not used with a.}
  \item \textbf{Incorrect}. Globalization and internationalization are related but do not always mean exactly the same thing. \emph{Globalization is a broad worldwide process, while internationalization often describes the deliberate international expansion of a particular organization.}
\end{enumerate}

\subsubsection{Collocation practice}
\begin{enumerate}
  \item \textbf{on.} The report examines the impact of globalization on rural communities. \emph{Use the impact of globalization on a group, place, or activity.}
  \item \textbf{accelerate.} New transportation systems helped accelerate globalization during the twentieth century. \emph{Accelerate globalization means cause the process to develop more rapidly.}
  \item \textbf{of.} The globalization of trade has connected producers with distant markets. \emph{Use the globalization of something.}
  \item \textbf{costs.} The course considers both the benefits and the costs of globalization. \emph{The costs of globalization refers to its negative consequences.}
\end{enumerate}

\vfill
\begin{center}
\small\color{Muted} Created and compiled by \textbf{Amir Kooshky}\\
\href{https://t.me/KooshkyTOEFL}{Telegram Channel}\quad\textbullet\quad
\href{https://www.instagram.com/kooshkytoefl}{Instagram}\quad\textbullet\quad
\href{https://t.me/KooshkyTOEFL\_pv}{Personal Telegram}
\end{center}
\end{document}
